% Spiegelmethode: Ladung vor geerdeter Ebene == Ladungspaar +Q/-Q (für y > 0)
% Gemeinsamer Stil für alle Abbildungen (TikZ -> SVG, siehe figures/build.mjs).
%
% Farben sind Platzhalter ("Sentinels"): build.mjs ersetzt sie im SVG durch
% CSS-Variablen, damit die Abbildungen im hellen und dunklen Modus passen.
% Andere Farben (auch schwarz/weiß) bricht der Build mit Fehler ab.
%
%   fg    Linien, Beschriftung          -> --dark
%   mu    Hilfslinien, Achsen, Maße      -> --gray
%   acc   das eine Konzept der Abbildung -> --secondary (Lime/Oliv)
%   blue  zweite Größe, negative Ladung  -> --fig-blue
%   warm  dritte Größe, positive Ladung  -> --fig-warm
%   bg    Seitenhintergrund (Aussparen)  -> --light
%
% Kein \clip verwenden: pdftocairo macht daraus Clip-Gruppen, in denen exakt
% waagerechte/senkrechte Linien verschwinden. Berechnete Kurven schneidet
% gen/*.py selbst am Rahmen ab.
\documentclass[tikz,border=3pt]{standalone}
\usepackage{amsmath,amssymb,bm}
% Text serifenlos wie die Seite, Mathe in Computer Modern wie KaTeX
\renewcommand{\familydefault}{\sfdefault}
\usetikzlibrary{arrows.meta,calc,decorations.markings,patterns,3d,positioning,angles,quotes}

\definecolor{fg}{RGB}{1,2,3}
\definecolor{mu}{RGB}{4,5,6}
\definecolor{acc}{RGB}{7,8,9}
\definecolor{blue}{RGB}{10,11,12}
\definecolor{warm}{RGB}{13,14,15}
\definecolor{bg}{RGB}{16,17,18}

% Vektoren wie auf der Website (KaTeX \mathbf)
\newcommand{\vv}[1]{\mathbf{#1}}

\tikzset{
  every picture/.style={fg, line width=0.8pt, line cap=round, line join=round,
    >={Stealth[length=5.5pt,width=4.4pt]}},
  every node/.style={font=\small},
  % Linientypen
  main/.style={line width=0.9pt},
  thin/.style={line width=0.5pt},
  aux/.style={mu, line width=0.5pt, dash pattern=on 2.5pt off 2pt},
  axis/.style={mu, line width=0.6pt, ->},
  vec/.style={line width=1.1pt, ->},
  field/.style={line width=0.7pt, postaction={decorate},
    decoration={markings, mark=at position #1 with {\arrow{Stealth[length=4.5pt,width=3.6pt]}}}},
  field/.default=0.5,
  % Flächen
  soft/.style={fill=#1, fill opacity=0.14},
  soft/.default=acc,
  % Ladungen
  qpos/.style={circle, fill=warm, inner sep=0pt, minimum size=9pt,
    path picture={\draw[bg, line width=0.9pt] ($(path picture bounding box.center)+(-2pt,0)$) -- +(4pt,0)
      ($(path picture bounding box.center)+(0,-2pt)$) -- +(0,4pt);}},
  qneg/.style={circle, fill=blue, inner sep=0pt, minimum size=9pt,
    path picture={\draw[bg, line width=0.9pt] ($(path picture bounding box.center)+(-2pt,0)$) -- +(4pt,0);}},
  dot/.style={circle, fill=fg, inner sep=0pt, minimum size=3pt},
  % Beschriftung mit Aussparung, damit Linien den Text nicht kreuzen
  lbl/.style={fill=bg, inner sep=1.5pt},
  % Icons: kräftigere Linien, keine Schrift
  icon/.style={line width=1.6pt, >={Stealth[length=6pt,width=5pt]}},
}

\begin{document}
\begin{tikzpicture}
  \def\a{1.1}
  % ---------- links: reale Anordnung ----------
  \begin{scope}
    \draw[field=0.45] plot coordinates {(0.136,1.131) (0.166,1.138) (0.195,1.144) (0.224,1.150) (0.254,1.157) (0.283,1.163) (0.313,1.169) (0.342,1.175) (0.371,1.180) (0.401,1.186) (0.430,1.191) (0.460,1.196) (0.490,1.201) (0.519,1.206) (0.549,1.210) (0.579,1.215) (0.608,1.219) (0.638,1.223) (0.668,1.226) (0.698,1.229) (0.727,1.232) (0.757,1.235) (0.787,1.238) (0.817,1.240) (0.847,1.241) (0.877,1.243) (0.907,1.244) (0.937,1.245) (0.967,1.246) (0.997,1.246) (1.027,1.246) (1.057,1.245) (1.087,1.244) (1.117,1.243) (1.147,1.241) (1.177,1.240) (1.207,1.237) (1.237,1.234) (1.267,1.231) (1.296,1.228) (1.326,1.224) (1.356,1.219) (1.385,1.215) (1.415,1.209) (1.444,1.204) (1.474,1.198) (1.503,1.191) (1.532,1.184) (1.561,1.177) (1.590,1.169) (1.619,1.161) (1.648,1.152) (1.676,1.143) (1.705,1.133) (1.733,1.123) (1.761,1.112) (1.789,1.101) (1.816,1.089) (1.844,1.077) (1.871,1.064) (1.898,1.051) (1.925,1.037) (1.951,1.023) (1.977,1.009) (2.003,0.994) (2.029,0.978) (2.054,0.962) (2.079,0.945) (2.104,0.928) (2.128,0.910) (2.152,0.892) (2.175,0.873) (2.198,0.854) (2.221,0.834) (2.243,0.814) (2.265,0.793) (2.286,0.772) (2.307,0.750)};
\draw[field=0.45] plot coordinates {(0.109,1.187) (0.133,1.206) (0.157,1.224) (0.180,1.243) (0.204,1.261) (0.228,1.280) (0.251,1.298) (0.275,1.316) (0.299,1.335) (0.323,1.353) (0.347,1.371) (0.371,1.388) (0.395,1.406) (0.419,1.424) (0.444,1.441) (0.468,1.459) (0.493,1.476) (0.517,1.493) (0.542,1.510) (0.567,1.527) (0.592,1.544) (0.616,1.561) (0.641,1.577) (0.667,1.594) (0.692,1.610) (0.717,1.626) (0.742,1.642) (0.768,1.658) (0.793,1.674) (0.819,1.689) (0.845,1.705) (0.871,1.720) (0.897,1.735) (0.923,1.750) (0.949,1.765) (0.975,1.780) (1.001,1.794) (1.028,1.808) (1.054,1.822) (1.081,1.836) (1.107,1.850) (1.134,1.864) (1.161,1.877) (1.188,1.890) (1.215,1.903) (1.242,1.916) (1.269,1.929) (1.296,1.942) (1.324,1.954) (1.351,1.966) (1.379,1.978) (1.406,1.990) (1.434,2.002) (1.461,2.013) (1.489,2.024) (1.517,2.035) (1.545,2.046) (1.573,2.057) (1.601,2.067) (1.630,2.078) (1.658,2.088) (1.686,2.098) (1.714,2.107) (1.743,2.117) (1.771,2.126) (1.800,2.135) (1.829,2.144) (1.857,2.153) (1.886,2.161) (1.915,2.169) (1.944,2.177) (1.973,2.185) (2.002,2.193) (2.031,2.200) (2.060,2.207) (2.089,2.214) (2.119,2.221) (2.148,2.228) (2.177,2.234) (2.207,2.240) (2.236,2.246) (2.265,2.252) (2.295,2.257) (2.305,2.259)};
\draw[field=0.45] plot coordinates {(0.061,1.226) (0.074,1.253) (0.087,1.280) (0.100,1.307) (0.113,1.334) (0.126,1.361) (0.140,1.388) (0.153,1.415) (0.166,1.442) (0.180,1.469) (0.193,1.495) (0.207,1.522) (0.220,1.549) (0.234,1.576) (0.247,1.602) (0.261,1.629) (0.275,1.656) (0.289,1.682) (0.303,1.709) (0.316,1.736) (0.330,1.762) (0.345,1.789) (0.359,1.815) (0.373,1.841) (0.387,1.868) (0.401,1.894) (0.416,1.920) (0.430,1.947) (0.445,1.973) (0.459,1.999) (0.474,2.025) (0.489,2.052) (0.504,2.078) (0.518,2.104) (0.533,2.130) (0.548,2.156) (0.563,2.182) (0.579,2.208) (0.594,2.233) (0.609,2.259) (0.624,2.285) (0.640,2.311) (0.655,2.337) (0.671,2.362) (0.686,2.388) (0.702,2.413) (0.718,2.439) (0.734,2.464) (0.749,2.490) (0.765,2.515) (0.781,2.541) (0.797,2.566) (0.814,2.591) (0.824,2.608)};
\draw[field=0.45] plot coordinates {(0.000,1.240) (0.000,1.270) (0.000,1.300) (0.000,1.330) (0.000,1.360) (0.000,1.390) (0.000,1.420) (0.000,1.450) (0.000,1.480) (0.000,1.510) (0.000,1.540) (0.000,1.570) (0.000,1.600) (0.000,1.630) (0.000,1.660) (0.000,1.690) (0.000,1.720) (0.000,1.750) (0.000,1.780) (0.000,1.810) (0.000,1.840) (0.000,1.870) (0.000,1.900) (0.000,1.930) (0.000,1.960) (0.000,1.990) (0.000,2.020) (0.000,2.050) (0.000,2.080) (0.000,2.110) (0.000,2.140) (0.000,2.170) (0.000,2.200) (0.000,2.230) (0.000,2.260) (0.000,2.290) (0.000,2.320) (0.000,2.350) (0.000,2.380) (0.000,2.410) (0.000,2.440) (0.000,2.470) (0.000,2.500) (0.000,2.530) (0.000,2.560) (0.000,2.590) (0.000,2.610)};
\draw[field=0.45] plot coordinates {(-0.061,1.226) (-0.074,1.253) (-0.087,1.280) (-0.100,1.307) (-0.113,1.334) (-0.126,1.361) (-0.140,1.388) (-0.153,1.415) (-0.166,1.442) (-0.180,1.469) (-0.193,1.495) (-0.207,1.522) (-0.220,1.549) (-0.234,1.576) (-0.247,1.602) (-0.261,1.629) (-0.275,1.656) (-0.289,1.682) (-0.303,1.709) (-0.316,1.736) (-0.330,1.762) (-0.345,1.789) (-0.359,1.815) (-0.373,1.841) (-0.387,1.868) (-0.401,1.894) (-0.416,1.920) (-0.430,1.947) (-0.445,1.973) (-0.459,1.999) (-0.474,2.025) (-0.489,2.052) (-0.504,2.078) (-0.518,2.104) (-0.533,2.130) (-0.548,2.156) (-0.563,2.182) (-0.579,2.208) (-0.594,2.233) (-0.609,2.259) (-0.624,2.285) (-0.640,2.311) (-0.655,2.337) (-0.671,2.362) (-0.686,2.388) (-0.702,2.413) (-0.718,2.439) (-0.734,2.464) (-0.749,2.490) (-0.765,2.515) (-0.781,2.541) (-0.797,2.566) (-0.814,2.591) (-0.824,2.608)};
\draw[field=0.45] plot coordinates {(-0.109,1.187) (-0.133,1.206) (-0.157,1.224) (-0.180,1.243) (-0.204,1.261) (-0.228,1.280) (-0.251,1.298) (-0.275,1.316) (-0.299,1.335) (-0.323,1.353) (-0.347,1.371) (-0.371,1.388) (-0.395,1.406) (-0.419,1.424) (-0.444,1.441) (-0.468,1.459) (-0.493,1.476) (-0.517,1.493) (-0.542,1.510) (-0.567,1.527) (-0.592,1.544) (-0.616,1.561) (-0.641,1.577) (-0.667,1.594) (-0.692,1.610) (-0.717,1.626) (-0.742,1.642) (-0.768,1.658) (-0.793,1.674) (-0.819,1.689) (-0.845,1.705) (-0.871,1.720) (-0.897,1.735) (-0.923,1.750) (-0.949,1.765) (-0.975,1.780) (-1.001,1.794) (-1.028,1.808) (-1.054,1.822) (-1.081,1.836) (-1.107,1.850) (-1.134,1.864) (-1.161,1.877) (-1.188,1.890) (-1.215,1.903) (-1.242,1.916) (-1.269,1.929) (-1.296,1.942) (-1.324,1.954) (-1.351,1.966) (-1.379,1.978) (-1.406,1.990) (-1.434,2.002) (-1.461,2.013) (-1.489,2.024) (-1.517,2.035) (-1.545,2.046) (-1.573,2.057) (-1.601,2.067) (-1.630,2.078) (-1.658,2.088) (-1.686,2.098) (-1.714,2.107) (-1.743,2.117) (-1.771,2.126) (-1.800,2.135) (-1.829,2.144) (-1.857,2.153) (-1.886,2.161) (-1.915,2.169) (-1.944,2.177) (-1.973,2.185) (-2.002,2.193) (-2.031,2.200) (-2.060,2.207) (-2.089,2.214) (-2.119,2.221) (-2.148,2.228) (-2.177,2.234) (-2.207,2.240) (-2.236,2.246) (-2.265,2.252) (-2.295,2.257) (-2.305,2.259)};
\draw[field=0.45] plot coordinates {(-0.136,1.131) (-0.166,1.138) (-0.195,1.144) (-0.224,1.150) (-0.254,1.157) (-0.283,1.163) (-0.313,1.169) (-0.342,1.175) (-0.371,1.180) (-0.401,1.186) (-0.430,1.191) (-0.460,1.196) (-0.490,1.201) (-0.519,1.206) (-0.549,1.210) (-0.579,1.215) (-0.608,1.219) (-0.638,1.223) (-0.668,1.226) (-0.698,1.229) (-0.727,1.232) (-0.757,1.235) (-0.787,1.238) (-0.817,1.240) (-0.847,1.241) (-0.877,1.243) (-0.907,1.244) (-0.937,1.245) (-0.967,1.246) (-0.997,1.246) (-1.027,1.246) (-1.057,1.245) (-1.087,1.244) (-1.117,1.243) (-1.147,1.241) (-1.177,1.240) (-1.207,1.237) (-1.237,1.234) (-1.267,1.231) (-1.296,1.228) (-1.326,1.224) (-1.356,1.219) (-1.385,1.215) (-1.415,1.209) (-1.444,1.204) (-1.474,1.198) (-1.503,1.191) (-1.532,1.184) (-1.561,1.177) (-1.590,1.169) (-1.619,1.161) (-1.648,1.152) (-1.676,1.143) (-1.705,1.133) (-1.733,1.123) (-1.761,1.112) (-1.789,1.101) (-1.816,1.089) (-1.844,1.077) (-1.871,1.064) (-1.898,1.051) (-1.925,1.037) (-1.951,1.023) (-1.977,1.009) (-2.003,0.994) (-2.029,0.978) (-2.054,0.962) (-2.079,0.945) (-2.104,0.928) (-2.128,0.910) (-2.152,0.892) (-2.175,0.873) (-2.198,0.854) (-2.221,0.834) (-2.243,0.814) (-2.265,0.793) (-2.286,0.772) (-2.307,0.750)};
\draw[field=0.45] plot coordinates {(-0.136,1.069) (-0.166,1.062) (-0.195,1.055) (-0.224,1.048) (-0.253,1.041) (-0.282,1.034) (-0.311,1.026) (-0.340,1.018) (-0.369,1.011) (-0.398,1.002) (-0.427,0.994) (-0.456,0.986) (-0.484,0.977) (-0.513,0.968) (-0.541,0.958) (-0.570,0.948) (-0.598,0.938) (-0.626,0.928) (-0.654,0.917) (-0.682,0.906) (-0.710,0.894) (-0.737,0.882) (-0.764,0.870) (-0.792,0.857) (-0.818,0.844) (-0.845,0.830) (-0.871,0.816) (-0.898,0.801) (-0.923,0.785) (-0.949,0.770) (-0.974,0.753) (-0.999,0.736) (-1.023,0.719) (-1.047,0.701) (-1.071,0.682) (-1.094,0.663) (-1.117,0.644) (-1.139,0.623) (-1.160,0.602) (-1.181,0.581) (-1.201,0.559) (-1.221,0.536) (-1.240,0.513) (-1.259,0.489) (-1.276,0.465) (-1.293,0.440) (-1.309,0.415) (-1.324,0.389) (-1.338,0.362) (-1.351,0.335) (-1.364,0.308) (-1.375,0.280) (-1.385,0.252) (-1.394,0.224) (-1.403,0.195) (-1.410,0.166) (-1.416,0.136) (-1.420,0.107) (-1.424,0.077) (-1.426,0.047) (-1.428,0.017) (-1.428,0.007)};
\draw[field=0.45] plot coordinates {(-0.109,1.013) (-0.133,0.994) (-0.156,0.975) (-0.179,0.956) (-0.203,0.937) (-0.226,0.918) (-0.249,0.898) (-0.271,0.879) (-0.294,0.859) (-0.317,0.840) (-0.339,0.820) (-0.361,0.799) (-0.383,0.779) (-0.405,0.758) (-0.426,0.737) (-0.447,0.716) (-0.468,0.694) (-0.489,0.673) (-0.509,0.650) (-0.529,0.628) (-0.548,0.605) (-0.567,0.582) (-0.586,0.558) (-0.604,0.534) (-0.621,0.510) (-0.638,0.485) (-0.654,0.459) (-0.669,0.434) (-0.684,0.407) (-0.698,0.381) (-0.711,0.354) (-0.723,0.327) (-0.735,0.299) (-0.745,0.271) (-0.755,0.242) (-0.764,0.214) (-0.771,0.185) (-0.778,0.155) (-0.783,0.126) (-0.787,0.096) (-0.790,0.066) (-0.792,0.036) (-0.793,0.006)};
\draw[field=0.45] plot coordinates {(-0.061,0.974) (-0.074,0.947) (-0.087,0.920) (-0.099,0.893) (-0.112,0.865) (-0.125,0.838) (-0.137,0.811) (-0.150,0.784) (-0.162,0.756) (-0.174,0.729) (-0.186,0.701) (-0.198,0.674) (-0.210,0.646) (-0.221,0.618) (-0.232,0.591) (-0.243,0.563) (-0.254,0.535) (-0.264,0.506) (-0.274,0.478) (-0.283,0.450) (-0.292,0.421) (-0.301,0.392) (-0.309,0.363) (-0.317,0.334) (-0.324,0.305) (-0.330,0.276) (-0.336,0.247) (-0.342,0.217) (-0.346,0.187) (-0.351,0.158) (-0.354,0.128) (-0.357,0.098) (-0.359,0.068) (-0.360,0.038) (-0.360,0.008)};
\draw[field=0.45] plot coordinates {(-0.000,0.960) (-0.000,0.930) (-0.000,0.900) (-0.000,0.870) (-0.000,0.840) (-0.000,0.810) (-0.000,0.780) (-0.000,0.750) (-0.000,0.720) (-0.000,0.690) (-0.000,0.660) (-0.000,0.630) (-0.000,0.600) (-0.000,0.570) (-0.000,0.540) (-0.000,0.510) (-0.000,0.480) (-0.000,0.450) (-0.000,0.420) (-0.000,0.390) (-0.000,0.360) (-0.000,0.330) (-0.000,0.300) (-0.000,0.270) (-0.000,0.240) (-0.000,0.210) (-0.000,0.180) (-0.000,0.150) (-0.000,0.120) (-0.000,0.090) (-0.000,0.060) (-0.000,0.030) (-0.000,0.010)};
\draw[field=0.45] plot coordinates {(0.061,0.974) (0.074,0.947) (0.087,0.920) (0.099,0.893) (0.112,0.865) (0.125,0.838) (0.137,0.811) (0.150,0.784) (0.162,0.756) (0.174,0.729) (0.186,0.701) (0.198,0.674) (0.210,0.646) (0.221,0.618) (0.232,0.591) (0.243,0.563) (0.254,0.535) (0.264,0.506) (0.274,0.478) (0.283,0.450) (0.292,0.421) (0.301,0.392) (0.309,0.363) (0.317,0.334) (0.324,0.305) (0.330,0.276) (0.336,0.247) (0.342,0.217) (0.346,0.187) (0.351,0.158) (0.354,0.128) (0.357,0.098) (0.359,0.068) (0.360,0.038) (0.360,0.008)};
\draw[field=0.45] plot coordinates {(0.109,1.013) (0.133,0.994) (0.156,0.975) (0.179,0.956) (0.203,0.937) (0.226,0.918) (0.249,0.898) (0.271,0.879) (0.294,0.859) (0.317,0.840) (0.339,0.820) (0.361,0.799) (0.383,0.779) (0.405,0.758) (0.426,0.737) (0.447,0.716) (0.468,0.694) (0.489,0.673) (0.509,0.650) (0.529,0.628) (0.548,0.605) (0.567,0.582) (0.586,0.558) (0.604,0.534) (0.621,0.510) (0.638,0.485) (0.654,0.459) (0.669,0.434) (0.684,0.407) (0.698,0.381) (0.711,0.354) (0.723,0.327) (0.735,0.299) (0.745,0.271) (0.755,0.242) (0.764,0.214) (0.771,0.185) (0.778,0.155) (0.783,0.126) (0.787,0.096) (0.790,0.066) (0.792,0.036) (0.793,0.006)};
\draw[field=0.45] plot coordinates {(0.136,1.069) (0.166,1.062) (0.195,1.055) (0.224,1.048) (0.253,1.041) (0.282,1.034) (0.311,1.026) (0.340,1.018) (0.369,1.011) (0.398,1.002) (0.427,0.994) (0.456,0.986) (0.484,0.977) (0.513,0.968) (0.541,0.958) (0.570,0.948) (0.598,0.938) (0.626,0.928) (0.654,0.917) (0.682,0.906) (0.710,0.894) (0.737,0.882) (0.764,0.870) (0.792,0.857) (0.818,0.844) (0.845,0.830) (0.871,0.816) (0.898,0.801) (0.923,0.785) (0.949,0.770) (0.974,0.753) (0.999,0.736) (1.023,0.719) (1.047,0.701) (1.071,0.682) (1.094,0.663) (1.117,0.644) (1.139,0.623) (1.160,0.602) (1.181,0.581) (1.201,0.559) (1.221,0.536) (1.240,0.513) (1.259,0.489) (1.276,0.465) (1.293,0.440) (1.309,0.415) (1.324,0.389) (1.338,0.362) (1.351,0.335) (1.364,0.308) (1.375,0.280) (1.385,0.252) (1.394,0.224) (1.403,0.195) (1.410,0.166) (1.416,0.136) (1.420,0.107) (1.424,0.077) (1.426,0.047) (1.428,0.017) (1.428,0.007)};

    % geerdete Ebene
    \draw[main] (-2.3,0) -- (2.3,0);
    % influenzierte Flächenladung sigma(x) ~ -a/(x^2+a^2)^{3/2}
    \fill[blue, fill opacity=0.35, domain=-2.3:2.3, samples=80]
      plot (\x, {-0.75*\a^3/((\x)^2+\a^2)^1.5}) -- (2.3,0) -- (-2.3,0) -- cycle;
    \draw[blue, thin, domain=-2.3:2.3, samples=80] plot (\x, {-0.75*\a^3/((\x)^2+\a^2)^1.5});
  \end{scope}
  \draw[main] (2.0,0) -- (2.0,-0.35) (1.8,-0.35) -- (2.2,-0.35);
  \draw[thin] (1.87,-0.43) -- (2.13,-0.43) (1.94,-0.51) -- (2.06,-0.51);
  \node[qpos] (Q) at (0,\a) {};
  \node[lbl, anchor=west] at (0.2,\a+0.22) {$+Q$};
  \draw[aux, |-|] (-1.95,0) -- node[lbl, font=\footnotesize, text=mu] {$a$} (-1.95,\a);
  \node[blue, font=\footnotesize, anchor=north] at (0.9,-0.72) {$\sigma(x)<0$};
  \node[mu, font=\footnotesize, anchor=west] at (2.22,-0.4) {$\phi=0$};
  \node[font=\small, anchor=south] at (0,2.65) {Leiter (geerdet)};

  % ---------- Äquivalenz ----------
  \node[acc, font=\Large] at (3.05,0.65) {$\Leftrightarrow$};
  \node[acc, font=\footnotesize, align=center] at (3.05,0.05) {gleich\\für $y>0$};

  % ---------- rechts: Ersatzanordnung ----------
  \begin{scope}[xshift=6.1cm]
    \input{data/spiegelmethode-bild}
    \draw[aux] (-2.3,0) -- (2.3,0);
    \node[qpos] at (0,\a) {};
    \node[lbl, anchor=west] at (0.2,\a+0.22) {$+Q$};
    \node[qneg] at (0,-\a) {};
    \node[anchor=west, text=mu] at (0.2,-\a-0.22) {$-Q$ (Bildladung)};
    \draw[aux, |-|] (-1.95,0) -- node[lbl, font=\footnotesize, text=mu] {$a$} (-1.95,-\a);
    \node[font=\small, anchor=south] at (0,2.65) {Ersatzproblem};
  \end{scope}
\end{tikzpicture}
\end{document}
